\chapter*{Introduction générale}
\addcontentsline{toc}{chapter}{Introduction générale}

La modernisation de l'administration publique passe aujourd'hui, dans de nombreux pays, par la numérisation de processus qui reposaient traditionnellement sur des démarches papier et une transmission hiérarchique manuelle des documents. La gestion des congés du personnel constitue un exemple particulièrement représentatif de cette problématique : elle concerne l'ensemble des agents d'une organisation, elle implique plusieurs niveaux de validation successifs, et elle doit produire, en fin de parcours, un document officiel dont l'authenticité ne peut être mise en doute.

Au sein du Ministère de l'Économie, des Finances et du Budget de la République de Guinée, ce processus reposait encore sur des méthodes essentiellement manuelles. Une demande de congé devait être formulée sur papier, transmise physiquement d'un service à l'autre, puis validée successivement par le chef de service, le directeur et le chef de cabinet, avant qu'une attestation ne soit délivrée à l'agent concerné. Cette organisation, bien que fonctionnelle, présentait plusieurs limites : des délais de traitement difficiles à maîtriser, une traçabilité incomplète des décisions prises à chaque étape, et surtout, aucune garantie technique permettant de vérifier qu'une attestation présentée par un agent n'avait pas été falsifiée.

C'est dans ce contexte que s'inscrit le présent stage, dont l'objectif était de concevoir et de développer une application capable de dématérialiser l'ensemble de ce processus, du dépôt de la demande par l'agent jusqu'à la délivrance d'une attestation authentifiable, en respectant fidèlement le circuit hiérarchique propre au ministère.

La démarche suivie a débuté par une analyse approfondie d'une solution libre déjà existante dans le domaine de la gestion des ressources humaines, afin d'identifier les mécanismes réutilisables et les limites à éviter, avant qu'un choix technologique définitif ne soit arrêté avec le ministère. Le développement s'est ensuite déroulé de manière itérative, chaque fonctionnalité étant accompagnée de tests automatisés destinés à garantir sa fiabilité avant son intégration au reste de l'application.

Ce rapport rend compte de l'ensemble de cette démarche. Il présente d'abord le contexte du stage et l'analyse de l'existant qui a orienté les choix initiaux du projet, avant d'exposer les besoins fonctionnels et non fonctionnels recensés. Il détaille ensuite l'architecture technique retenue et les diagrammes de conception associés, puis décrit les principales réalisations du projet. Une attention particulière est portée aux difficultés effectivement rencontrées au cours du développement et aux solutions qui y ont été apportées, avant de conclure sur le bilan de ce stage et les perspectives d'évolution de l'application.
